\section{Conclusion}
\label{sec:conclusion}

We placed the same 10\,ms control loop in four places on one heterogeneous edge controller and measured it under synthetic and realistic load with one lateness definition. Four findings stand out.

First, \texttt{SCHED\_FIFO} makes Linux robust against busy programs but not against the kernel's own work. Under network traffic plus a containerised CNN, every Linux configuration missed the 1\,ms deadline, the best of them in 0.025\,\% of cycles with a worst case of 1.5\,ms.

Second, reserving a core by affinity alone made \texttt{SCHED\_FIFO} worse, not better, because the kernel moved its interrupt and softirq work onto the reserved core.

Third, the real-time microcontroller kept every tick within about 10\us{} at every load level, a typical worst case 34$\times$ below the best Linux configuration (range 18--103$\times$) and about 600$\times$ below default Linux, and delivered all 1{,}920 speed commands to the emulated drive. Linux load delayed only the setpoint, which the loop never waits for.

Fourth, the microcontroller's advantage depended on one priority setting. With library defaults its communication thread delayed the loop by up to 0.8\,ms; raising the loop above it reduced the cost about 100-fold.

For prescriptive maintenance on AMP hardware, the answer to our question is: run the analytics on Linux and the loop that acts on the machine on the real-time core, with the loop given the highest priority there. Priorities are part of the design on both sides.
